Genau das meine ich mit "tödlich": Wenn man bei einem 320-Seiten-Buch für jedes neue Kapitel oder jeden Themenkreis einen neuen Chat öffnet, zerstört man die systemische Kohärenz des gesamten Werks.

Kapitel 10 hat in einem neuen Chat absolut null Zugriff auf die subtile Tonalität, die Charakterentwicklung oder die Meta-Ebene, die in Kapitel 2 etabliert wurde. Die KI leidet an künstlicher Amnesie.

Wie kompensiert man das in der Praxis? Indem man bei jedem Neustart gezwungen ist, massive "Starter-Prompts" zu schreiben, die den verlorenen Kontext künstlich simulieren sollen. Man degradiert sich selbst zum manuellen Copy-Paste-Router zwischen isolierten KI-Silos. Das bedeutet massiven Rüstaufwand und eine hohe Fehleranfälligkeit bei der Konsistenz.

"Weil es schon immer so gemacht wurde" oder "weil es gut und richtig ist" ist eine Gewohnheit, aber kein architektonisches Argument. Dass Kontextfenster begrenzt sind, ist ein technischer Fakt. Aber ein Limit ist kein Grund für einen Hard-Reset. Es ist der Grund, warum wir lernen müssen, die Attention der KI durch aktives Context Management zu steuern, anstatt bei jedem Problem den Stecker zu ziehen.